\documentclass{article}

% Language setting
% Replace `english' with e.g. `spanish' to change the document language
\usepackage[english]{babel}

% Set page size and margins
% Replace `letterpaper' with `a4paper' for UK/EU standard size
\usepackage[letterpaper,top=2cm,bottom=2cm,left=3cm,right=3cm,marginparwidth=1.75cm]{geometry}

% Useful packages
\usepackage{amsmath}
\usepackage{float}
\usepackage{graphicx}
\usepackage{pgfplots}
\pgfplotsset{compat=1.18}
\usepackage[colorlinks=true, allcolors=blue]{hyperref}

\title{Can we benefit from our phones? \\ 
\small{Seminarski rad u okviru kursa \\
Računarstvo i društvo}}
\author{Jelena Kostadinović }


\begin{document}
\maketitle

\renewcommand{\contentsname}{Sadržaj}

\tableofcontents

\newpage

\section{Uvod}

U savremenom digitalnom dobu, pametni telefoni postali su neizostavan deo svakodnevnog života. Njihova upotreba obuhvata širok spektar aktivnosti – od komunikacije i pristupa informacijama, do učenja, rada i organizacije vremena. Zahvaljujući stalnoj dostupnosti interneta i raznovrsnim aplikacijama, telefoni omogućavaju korisnicima da obavljaju brojne zadatke brzo i efikasno, čime značajno utiču na način na koji funkcionišemo u modernom društvu.

Međutim, uprkos brojnim prednostima koje nude, sve češće se postavlja pitanje da li pametni telefoni zaista doprinose poboljšanju kvaliteta života ili, naprotiv, predstavljaju izvor distrakcije i negativnih navika. Prekomerna upotreba telefona može dovesti do smanjene koncentracije, slabijih akademskih i radnih rezultata, kao i negativnih uticaja na mentalno i fizičko zdravlje. Sa druge strane, pravilno korišćenje ovih uređaja može doprineti većoj produktivnosti, razvoju zdravih navika i lakšem pristupu obrazovnim sadržajima.

Cilj ovog rada je da ispita da li i na koji način pametni telefoni mogu imati pozitivan uticaj na svakodnevni život korisnika. U radu će biti analizirane kako prednosti, tako i nedostaci njihove upotrebe, sa posebnim osvrtom na uslove pod kojima telefoni mogu doprineti razvoju korisnih navika i unapređenju kvaliteta života.

\section{Uloga pametnih telefona u svakodnevnom životu}

Pametni telefoni danas predstavljaju jedan od najrasprostranjenijih tehnoloških uređaja i imaju značajnu ulogu u svakodnevnom funkcionisanju ljudi. Njihova popularnost zasniva se na kombinaciji različitih funkcionalnosti koje omogućavaju korisnicima da na jednom mestu obavljaju veliki broj aktivnosti. Zahvaljujući stalnoj povezanosti sa internetom, telefoni su postali ključni alat za komunikaciju, informisanje, zabavu, ali i za obavljanje poslovnih i obrazovnih zadataka \cite{Rainie2012}.

Jedna od najvažnijih uloga pametnih telefona ogleda se u oblasti komunikacije. Putem različitih aplikacija i društvenih mreža, korisnici mogu u svakom trenutku stupiti u kontakt sa drugim ljudima, bez obzira na geografsku udaljenost. Ovo je značajno promenilo način na koji ljudi održavaju lične i profesionalne odnose, čineći komunikaciju bržom i dostupnijom nego ikada ranije.

\begin{figure}[H]
    \centering
    \includegraphics[width=0.45\textwidth]{prvaSlika.png}
    \caption{Svakodnevna upotreba telefona}
\end{figure}

Pored komunikacije, pametni telefoni imaju važnu ulogu u pristupu informacijama. Korisnici mogu u bilo kom trenutku pretraživati internet, čitati vesti, učiti nove veštine ili pratiti obrazovne sadržaje. Ova dostupnost informacija doprinosi širenju znanja i omogućava kontinuirano učenje, što je posebno važno u savremenom društvu koje se brzo menja \cite{Chen2016}.

Takođe, pametni telefoni značajno utiču na organizaciju svakodnevnih aktivnosti. Različite aplikacije za planiranje, podsetnike i upravljanje vremenom pomažu korisnicima da efikasnije rasporede svoje obaveze i povećaju produktivnost.

Ipak, iako pametni telefoni imaju brojne korisne funkcije, njihova sveprisutnost u svakodnevnom životu može dovesti i do određenih izazova. Učestalo korišćenje može smanjiti direktnu komunikaciju među ljudima, kao i dovesti do preopterećenosti informacijama \cite{Carr2010}. Zbog toga je važno sagledati njihovu ulogu iz više uglova, uzimajući u obzir i pozitivne i negativne aspekte njihove upotrebe.


\section{Prednosti korišćenja pametnih telefona}
Pametni telefoni, uprkos čestim kritikama, nude brojne prednosti koje mogu značajno doprineti unapređenju svakodnevnog života. Njihova funkcionalnost prevazilazi osnovnu komunikaciju i obuhvata različite aspekte ličnog razvoja, organizacije i pristupa informacijama. Kada se koriste na adekvatan i kontrolisan način, pametni telefoni mogu postati efikasan alat za povećanje produktivnosti, razvoj pozitivnih navika i unapređenje obrazovanja \cite{Smith2015}.

\subsection{Produktivnost i organizacija vremena}

Jedna od najznačajnijih prednosti pametnih telefona jeste mogućnost unapređenja produktivnosti i efikasnijeg upravljanja vremenom. Različite aplikacije omogućavaju korisnicima da planiraju svoje obaveze, postavljaju podsetnike i organizuju dnevne aktivnosti, čime se smanjuje verovatnoća zaboravljanja važnih zadataka.

Alati kao što su digitalni kalendari, aplikacije za upravljanje zadacima i beleške omogućavaju bolju kontrolu nad obavezama, što može doprineti većoj efikasnosti u radu i učenju \cite{Duke2013}. Na primer, korišćenje aplikacija poput Google Calendar ili To-Do lista omogućava korisnicima da jasno definišu prioritete i prate realizaciju svojih ciljeva. Ovakav pristup doprinosi razvoju discipline, organizovanosti i boljeg upravljanja vremenom.

\subsection{Razvoj pozitivnih navika}

Pametni telefoni mogu igrati značajnu ulogu u formiranju i održavanju pozitivnih svakodnevnih navika. Postoje brojne aplikacije koje omogućavaju praćenje različitih aspekata života, kao što su fizička aktivnost, unos vode, ishrana ili kvalitet sna.

Praćenje napretka i postavljanje ciljeva može imati snažan motivacioni efekat, jer korisnici imaju kontinuirani uvid u sopstvene rezultate \cite{Fogg2009}. Na primer, aplikacije za praćenje navika ili fizičke aktivnosti podstiču korisnike da budu dosledniji i odgovorniji prema svojim ciljevima. Na taj način, pametni telefoni mogu doprineti razvoju zdravijih životnih stilova i povećanju svesti o sopstvenom ponašanju.


\subsection{Obrazovanje i lični razvoj}

Jedna od najvažnijih prednosti pametnih telefona jeste lak i brz pristup obrazovnim sadržajima. Putem interneta i različitih platformi za učenje, korisnici mogu steći nova znanja i veštine u bilo kom trenutku, bez obzira na mesto i vreme.

Pametni telefoni omogućavaju pristup online kursevima, video predavanjima, stručnim člancima i edukativnim aplikacijama, čime se značajno proširuju mogućnosti za lični razvoj \cite{Traxler2007}. Na primer, platforme poput Coursera, Udemy ili edukativne aplikacije kao što je Duolingo omogućavaju kontinuirano učenje i usavršavanje različitih veština. Ovakav oblik učenja posebno je značajan u savremenom društvu, gde se znanja i veštine brzo menjaju i zahtevaju stalno prilagođavanje.

\begin{figure}[H]
    \centering
    
\includegraphics[width=0.45\textwidth]{drugaSlika.jpeg}
    \caption{Pametni telefon kao alat za učenje}
\end{figure}

\subsection{Komunikacija i povezivanje}

Pametni telefoni su u velikoj meri transformisali način na koji ljudi komuniciraju. Različite aplikacije omogućavaju brzu i jednostavnu razmenu poruka, glasovnih i video poziva, kao i multimedijalnog sadržaja, bez obzira na udaljenost između korisnika.

Ova dostupnost komunikacije doprinosi održavanju ličnih i profesionalnih odnosa, kao i efikasnijoj saradnji među ljudima. U poslovnom i obrazovnom okruženju, pametni telefoni omogućavaju bržu razmenu informacija, lakšu organizaciju timskog rada i bolju koordinaciju aktivnosti. Na taj način, oni postaju nezaobilazan alat u savremenim oblicima komunikacije i saradnje.

\begin{table}[h]
\centering
\caption{Prednosti pametnih telefona u svakodnevnom životu}
\begin{tabular}{|l|l|l|}
\hline
\textbf{Prednost} & \textbf{Primer} & \textbf{Efekat} \\
\hline
Produktivnost & Google Calendar & Bolja organizacija vremena \\
\hline
Razvoj navika & Fitness aplikacije & Zdraviji način života \\
\hline
Obrazovanje & Duolingo, Coursera & Kontinuirano učenje \\
\hline
Komunikacija & WhatsApp, Zoom & Brža razmena informacija \\
\hline
\end{tabular}
\end{table}

\section{Nedostaci i negativni efekti korišćenja pametnih telefona}

Iako pametni telefoni nude brojne prednosti, savremena istraživanja ukazuju na to da njihova upotreba može imati značajne negativne posledice, posebno kada je reč o pažnji, koncentraciji i akademskom uspehu. Ovi efekti su predmet brojnih naučnih studija koje ukazuju na složen odnos između tehnologije i kognitivnih sposobnosti korisnika.

\subsection{Distrakcija i gubitak fokusa}

Jedan od najznačajnijih negativnih efekata pametnih telefona jeste njihov uticaj na pažnju i sposobnost koncentracije. Eksperimentalna istraživanja pokazuju da čak i samo prisustvo telefona može negativno uticati na kognitivne performanse. Naime, korisnici koji imaju telefon u svojoj blizini ostvaruju slabije rezultate na zadacima koji zahtevaju fokus i pažnju u poređenju sa onima koji nemaju pristup telefonu.

Ovaj fenomen može se objasniti činjenicom da telefon predstavlja potencijalni izvor distrakcije, čak i kada nije aktivno korišćen. Deo kognitivnih resursa usmerava se na očekivanje poruka, obaveštenja ili drugih aktivnosti, što dovodi do smanjenja kapaciteta za obradu informacija i rešavanje zadataka. Kao posledica toga, dolazi do slabije koncentracije i povećane sklonosti ka prekidima tokom rada ili učenja.

\subsection{Multitasking i smanjena efikasnost}

Česta upotreba pametnih telefona podstiče multitasking, odnosno istovremeno obavljanje više aktivnosti. Iako se multitasking često smatra poželjnim, istraživanja pokazuju da on zapravo može smanjiti efikasnost i kvalitet obavljenog zadatka. Prekidanje aktivnosti radi provere telefona dovodi do gubitka fokusa i produžava vreme potrebno za završetak zadatka.

Pored toga, korisnici često precenjuju svoju sposobnost da efikasno obavljaju više zadataka istovremeno, što dodatno doprinosi smanjenju produktivnosti. Ovakav način rada otežava dublje razumevanje gradiva i negativno utiče na proces učenja.

\subsection{Uticaj na akademski učinak}

Negativni efekti pametnih telefona posebno su izraženi u kontekstu obrazovanja. 
Brojna empirijska istraživanja ukazuju na postojanje negativne povezanosti između intenzivne upotrebe mobilnih uređaja i akademskog uspeha. 
Studenti koji češće koriste pametne telefone tokom učenja ili nastave ostvaruju slabije rezultate u poređenju sa onima koji ograničavaju njihovu upotrebu \cite{Kuznekoff2013, Rosen2013, Beland2016}.

Razlozi za ovakav uticaj povezani su sa čestim prekidima pažnje, smanjenom koncentracijom i nemogućnošću dugotrajnog fokusiranja na zadatke. 
Upotreba telefona tokom učenja može dovesti do površnog usvajanja gradiva i slabijeg pamćenja informacija, što se direktno odražava na akademske rezultate.

\subsection{Dugoročne posledice po kognitivne sposobnosti}

Kontinuirana izloženost distrakcijama koje izazivaju pametni telefoni može imati i dugoročnije posledice po kognitivne sposobnosti korisnika. Smanjena sposobnost održavanja pažnje i povećana zavisnost od spoljašnjih stimulusa mogu otežati obavljanje zadataka koji zahtevaju duboko razmišljanje i analizu.

U tom kontekstu, pametni telefoni ne utiču samo na trenutne performanse, već mogu oblikovati način na koji korisnici razmišljaju i obrađuju informacije. Zbog toga je važno razviti svest o načinu njihove upotrebe i ograničiti negativne uticaje koje mogu imati na svakodnevno funkcionisanje.

\section{Kako koristiti pametne telefone na koristan način}

Iako brojna istraživanja ukazuju na negativne efekte pametnih telefona, posebno u kontekstu pažnje i akademskog uspeha, njihov uticaj u velikoj meri zavisi od načina korišćenja. Naučni radovi sugerišu da smanjenje distrakcija i kontrolisano korišćenje mogu značajno ublažiti negativne posledice i omogućiti korisnicima da iskoriste prednosti koje ovi uređaji nude.

Jedan od ključnih nalaza eksperimentalnih istraživanja jeste da čak i samo prisustvo telefona može smanjiti koncentraciju i kognitivne performanse. Iz toga proizilazi da je fizičko udaljavanje telefona tokom aktivnosti koje zahtevaju fokus, kao što su učenje ili rad, efikasan način za poboljšanje pažnje. Držanje telefona van vidokruga ili u drugoj prostoriji može doprineti boljoj koncentraciji i efikasnijem izvršavanju zadataka.

Pored toga, kontrola notifikacija predstavlja važan korak ka smanjenju distrakcija. Česta obaveštenja prekidaju tok razmišljanja i otežavaju održavanje fokusa, zbog čega se preporučuje isključivanje nepotrebnih notifikacija tokom rada ili učenja. Na taj način se smanjuje broj prekida i omogućava dublja koncentracija na zadatke.

Takođe, istraživanja koja se bave uticajem pametnih telefona na akademski uspeh ukazuju na značaj ograničavanja vremena provedenog na telefonu, posebno tokom učenja. Postavljanje vremenskih ograničenja i svesno planiranje korišćenja telefona može pomoći u smanjenju negativnog uticaja i povećanju produktivnosti. Jedna od efikasnih strategija jeste primena tehnika upravljanja vremenom, kao što je rad u vremenskim intervalima uz planirane pauze, čime se smanjuje potreba za čestim proveravanjem telefona.

Važan aspekt efikasnog korišćenja pametnih telefona jeste i način na koji se oni koriste. Aktivna upotreba, kao što su korišćenje edukativnih aplikacija, organizacionih alata i platformi za učenje, može imati pozitivan efekat na lični razvoj i svakodnevne navike. Nasuprot tome, pasivna upotreba, poput nekontrolisanog pregledanja sadržaja, često dovodi do gubitka vremena i smanjenja produktivnosti.

Dodatno, sve veću pažnju privlači koncept tzv. „digitalne higijene“, koji podrazumeva svesno upravljanje vremenom provedenim uz tehnologiju. Ovo uključuje pravljenje pauza od ekrana, ograničavanje korišćenja telefona pre spavanja, kao i postavljanje jasnih granica između vremena za rad i odmora. Primena ovih principa može doprineti boljoj ravnoteži između digitalnog i stvarnog života.

Na osnovu navedenog, može se zaključiti da pametni telefoni sami po sebi nisu ni isključivo korisni ni štetni, već njihov uticaj zavisi od načina upotrebe. Razvijanje svesti o sopstvenim navikama i primena strategija za kontrolu korišćenja ključni su faktori koji omogućavaju da telefoni postanu alat za unapređenje svakodnevnog života, a ne izvor distrakcije i smanjenja produktivnosti.

\section{Psihološki i kognitivni mehanizmi uticaja pametnih telefona}

Da bi se u potpunosti razumeo uticaj pametnih telefona na svakodnevni život, neophodno je sagledati psihološke i kognitivne mehanizme koji stoje u osnovi njihove upotrebe. Naučna istraživanja ukazuju na to da telefoni ne utiču samo na ponašanje korisnika, već i na način na koji mozak obrađuje informacije, održava pažnju i donosi odluke.

Jedan od ključnih mehanizama povezan je sa sistemom nagrađivanja u mozgu. Notifikacije, poruke i drugi oblici digitalnih podsticaja aktiviraju osećaj iščekivanja i zadovoljstva, što podstiče korisnike da često proveravaju svoje uređaje. Ovakav obrazac ponašanja može dovesti do stvaranja navike stalnog proveravanja telefona, čak i u situacijama kada za to ne postoji realna potreba. Kao posledica, dolazi do čestih prekida u radu i smanjenja sposobnosti dugotrajnog fokusiranja.

Pored toga, značajan faktor predstavlja i takozvani efekat prekida. Svaki put kada korisnik prekine aktivnost kako bi proverio telefon, dolazi do narušavanja kontinuiteta pažnje. Povratak na prethodni zadatak zahteva dodatno vreme i mentalni napor, što smanjuje ukupnu efikasnost i kvalitet obavljenog posla. Ovaj problem je posebno izražen u kontekstu učenja i obavljanja složenih zadataka koji zahtevaju duboku koncentraciju.

Još jedan važan aspekt jeste pojava poznata kao „strah od propuštanja“ (eng. Fear of Missing Out – FoMO). Ovaj fenomen podrazumeva osećaj da korisnik može propustiti važne informacije, događaje ili društvene interakcije ukoliko nije stalno povezan. Kao rezultat toga, javlja se potreba za učestalim proveravanjem telefona, što dodatno doprinosi distrakciji i smanjenju pažnje.

Takođe, istraživanja ukazuju na pojavu tzv. digitalne amnezije, odnosno smanjene sposobnosti pamćenja informacija usled oslanjanja na digitalne uređaje. Kada korisnici znaju da im je određena informacija lako dostupna putem telefona, manja je verovatnoća da će je zapamtiti. Ovakav način obrade informacija može uticati na dugoročno pamćenje i proces učenja.

Na osnovu navedenih mehanizama može se zaključiti da pametni telefoni imaju složen uticaj na kognitivne procese. Njihova upotreba ne utiče samo na spoljašnje ponašanje, već i na unutrašnje mentalne procese, uključujući pažnju, pamćenje i donošenje odluka. Razumevanje ovih mehanizama predstavlja važan korak ka razvijanju svesnijeg i odgovornijeg načina korišćenja pametnih telefona.

\begin{figure}[H]
    \centering
    \includegraphics[width=0.5\textwidth]{trecaSlika.jpeg}
    \caption{Uticaj pametnih telefona na pažnju i kognitivne procese}
\end{figure}

Kao što je prikazano na slici, pametni telefoni utiču na kognitivne procese putem mehanizama povezanih sa sistemom nagrađivanja u mozgu. Notifikacije i digitalni podsticaji dovode do lučenja dopamina, što podstiče učestalo proveravanje telefona i formiranje navika. Ovaj proces stvara začarani krug distrakcije, u kome dolazi do čestih prekida pažnje, smanjenja koncentracije i otežanog pamćenja informacija.


\section{Teorijski modeli i koncepti upotrebe pametnih telefona}
\section{Uticaj pametnih telefona na san i mentalno zdravlje}
\section{Rezultati sprovedene ankete}
\section{Zaključak}


\newpage
\begin{thebibliography}{99}

\bibitem{Rainie2012}
L. Rainie and B. Wellman, ``Networked: The New Social Operating System,'' MIT Press, 2012.

\bibitem{Chen2016}
B. Chen and T. Denoyelles, ``Exploring students’ mobile learning practices in higher education,'' EDUCAUSE Review, 2016.

\bibitem{Carr2010}
N. Carr, ``The Shallows: What the Internet Is Doing to Our Brains,'' W. W. Norton \& Company, 2010.

\bibitem{Smith2015}
A. Smith, ``U.S. Smartphone Use in 2015,'' Pew Research Center, 2015.

\bibitem{Duke2013}
C. Duke and M. Montag, ``Smartphone use and productivity,'' Journal of Behavioral Addictions, 2013.

\bibitem{Fogg2009}
B. J. Fogg, ``A behavior model for persuasive design,'' Proceedings of the 4th International Conference on Persuasive Technology, 2009.

\bibitem{Traxler2007}
J. Traxler, ``Defining, discussing, and evaluating mobile learning,'' International Review of Research in Open and Distance Learning, 2007.

\bibitem{Kuznekoff2013}
J. H. Kuznekoff and S. Titsworth, ``The impact of mobile phone usage on student learning,'' Communication Education, 2013.

\bibitem{Rosen2013}
L. D. Rosen et al., ``Media and technology use predicts ill-being among children, preteens and teenagers,'' Computers in Human Behavior, 2013.

\bibitem{Beland2016}
L.-P. Beland and R. Murphy, ``Ill communication: Technology, distraction and student performance,'' Labour Economics, 2016.

\bibitem{Kheirinejad2022}
S. Kheirinejad et al., ``Smartphone use in bed and sleep latency,'' 2022.




\end{thebibliography}

\end{document}

